\begin{textsample}{POS Dim 1 – rr – Score 36.00 – rr/rr\_inf\_e\_ef\_5.txt}  \label{ex:f1_pos_265}
2.3 CURRÍCULO : A BNCC E OS DESAFIOS DE CONSTRUÇÃO DE \textbf{UMA} EDUCAÇÃO QUE ACOLHA A DIVERSIDADE E IDENTIDADE CULTURAL Localizado no extremo norte do Brasil, Roraima é \textbf{um} Estado com, aproximadamente, 522.6 mil habitantes ( BRASIL, 2017 ) e características muito diferenciadas em relação ao cenário nacional.
Com \textbf{uma} população indígena \textbf{que} representa, não apenas a maioria populacional, mas \textbf{que} também é muito diversa, envolvendo dez diferentes Povos, com línguas, costumes, tradições e ocupação territorial próprio.
O estado também faz \textbf{fronteira} com a Guiana e Venezuela, o \textbf{que} propicia a vinda de estrangeiros com a perspectiva de firmarem residência em seus 15 municípios.
A constituição populacional é fortemente marcada por movimentos migratórios de diversas regiões do país, mais acentuadamente de Estados como Paraná, Rio Grande do Sul, Maranhão e Pará.
Segundo o Censo Escolar ( BRASIL, 2017 ), A Rede Estadual de Educação de Roraima conta atualmente com 381 escolas, dentre as quais, 258 são escolas indígenas, distribuídas entre as 32 Terras Indígenas homologadas, \textbf{que} representa 46,2 \% do território do Estado, com algumas \textbf{que} já estão na zona urbana da capital Boa Vista.
Outro aspecto importante a ser destacado, é a migração maciça de venezuelanos \textbf{que}, de acordo com os dados da Secretaria Estadual de Educação, entre 2017 e 2018, gerou \textbf{uma} demanda de atendimento de mais de 1300 alunos nas classes de Ensino Fundamental II e Médio, e o número crescente de alunos a serem atendidos no âmbito da Educação Especial, \textbf{que} representa cerca de 1,7 \% do total de matrículas.
Esta breve descrição faz -se necessária para \textbf{que} tenhamos contornos mais claros sobre as demandas e desafios das Redes de Ensino, no \textbf{que} tange à responsabilidade do poder público em atender e incluir a diversidade \textbf{que} caracteriza suas unidades escolares, bem como dos diversos elementos a serem considerados nos processos de construção de seus Currículos e Projetos Político Pedagógicos escolares.
Temos \textbf{vivido}, numa escala mundial, movimentos \textbf{que} põem em debate aspectos e elementos importantes do campo da educação e suas necessidades de atendimento da diversidade, as políticas públicas, além dos papéis e significados da escola e do currículo no cotidiano, na e para a constituição \textbf{identitária} destes \textbf{sujeitos}.
Ao observarmos a escola em sua historicidade, percebemos \textbf{que} a diversidade sempre existiu.
A questão \textbf{central} é \textbf{que} \textbf{hoje}, este público tem feito exigências, reivindicações sobre seus direitos, cobrando \textbf{uma} nova perspectiva de ensino e formação escolar, de modo \textbf{que} \textbf{vejam} a pluralidade e o respeito às diferenças, a diversidade cultural, como valores positivos e isso tem \textbf{que} compor a pauta de discussões e construção dos currículos escolares.
Por esse motivo, a questão não pode ser refletida apenas em suas nuances \textbf{didático-pedagógicas}, voltadas para como e o \textbf{que} ensinar das línguas e/ou sobre as diversas culturas, ou seja, a interculturalidade nas escolas.
Corre -se o risco, com essa perspectiva, do apagamento dos elementos socioculturais e linguísticos como \textbf{constitutivas} dos indivíduos e determinantes aos modos de relação e organização social.
Os discursos e movimentos sociais têm apontado para \textbf{um} desejo de igualdade concreto, de outra ética educacional, \textbf{que} não pode ficar fora de \textbf{uma} pauta fundamental de reflexões e mudanças sociais.
Na visão de Moreira \& Candau, ( 2003, p. 160 ) : > A escola é, sem dúvidas, \textbf{uma} instituição cultural.
Portanto as relações entre escola e cultura não podem ser concebidas como entre dois polos independentes, mas sim como universos entrelaçados, como \textbf{uma} teia tecida no cotidiano com fios e nós profundamente articulados.
Neste cenário, a Base Nacional Curricular Comum ( BNCC ) precisa ser pensada e compreendida no \textbf{que} deseja encaminhar em sua proposta e princípios.
Em primeira \textbf{instância}, é \textbf{preciso} reafirmar a cada instante \textbf{que} a BNCC não tem como vocação prescrever, tão pouco reformar os currículos, ela vem para apontar caminhos, propostas curriculares a serem construídas por cada escola, considerando os direitos de aprendizagem \textbf{que} cada aluno brasileiro tem, tornando -se \textbf{um} documento de referência para a educação escolar, o \textbf{que} envolve as áreas de conhecimento, os componentes curriculares, mas também \textbf{um} movimento inevitável de reconfiguração das práticas de ensino, dos valores e da própria vida escolar, nos modos de lidar com a complexidade \textbf{que} lhe é \textbf{inerente}.
A BNCC aponta para as competências e habilidades a serem trabalhadas, determina \textbf{que} as redes de ensino orientem suas unidades para o exercício da autonomia na construção de seus projetos educacionais, mostrando \textbf{que} precisamos empreender propostas \textbf{que} reconheçam as diferenças e a interculturalidade como questões essenciais em qualquer escola \textbf{que} se proponha atender de maneira responsiva aos seus alunos.
O \textbf{que} é \textbf{um} grande desafio a ser vencido, \textbf{visto} \textbf{que} a educação não é \textbf{neutra}, \textbf{que} responde a interesses político ideológicos.
As reformas educacionais latino-americanas, marcadas por evidente diferenciação étnico-cultural de suas populações, declaram o propósito de superar os modelos colonizadores, \textbf{que} buscam \textbf{uma} adaptação unidirecional dos estudantes para modelos culturais homogêneos, e têm procurando substituí -los por modelos multiculturais para melhorar atenção educacional à diversidade ; no entanto, estes novos modelos propõem currículos culturalmente diferenciados e terminam se tornando excludentes, pois encapsulam a diversidade cultural no aluno e sem estabelecer as ideias de diálogo intercultural \textbf{que} os sustentam ( IBÁÑEZ-SALGADO \textbf{et} al, 2012 ).
Segundo Candau ( 2014, p. 28 ) : > [... ] o currículo deve ser intercultural, capaz de : promover deliberada \textbf{inter-relação} entre os diferentes \textbf{sujeitos} e grupos socioculturais presentes em determinada sociedade ; nesse sentido, essa posição se situa em confronto com todas as visões diferencialistas, assim como com as perspectivas assimilacionistas ; por outro \textbf{lado}, rompe com \textbf{uma} visão essencialista das culturas e das identidades culturais ; concebe as culturas em contínuo processo de construção, desestabilização e reconstrução ; está constituída pela afirmação de \textbf{que} nas sociedades em \textbf{que} \textbf{vivemos} os processos de hibridização cultural são intensos e mobilizadores de construção permanente, o \textbf{que} \textbf{supõe} \textbf{que} as culturas não são puras, nem estáticas ; tem presente os mecanismos de poder \textbf{que} \textbf{permeiam} as relações românticas, estão \textbf{atravessadas} por questões de poder e marcadas pelo preconceito e discriminação de determinados grupos socioculturais.
A interculturalidade, no cotidiano escolar, precisa considerar a vida dos alunos, as contradições presentes em suas realidades, as relações entre as diferentes sociedades e conhecimentos, se mantendo integradas as suas raízes e, ao mesmo tempo, conectadas ao global, trabalhando com tudo \textbf{que} estas duas dimensões oferecem.
A escola está chamada a ser, nos próximos anos, mais do \textbf{que} \textbf{um} \textbf{lócus} de apropriação do conhecimento socialmente relevante, o científico, \textbf{um} espaço de diálogo entre diferentes saberes – científico, social, escolar, etc. – e linguagens.
De análise crítica, \textbf{estímulo} ao exercício da capacidade reflexiva e de \textbf{uma} visão plural e histórica do conhecimento, da ciência, da tecnologia e das diferentes linguagens.
É no cruzamento, na interação, no reconhecimento da dimensão histórica e social \textbf{que} a escola está chamada a se situar ( CANDAU, 2008, p. 14 ).
No caso de escolas indígenas, por exemplo, os valores, os costumes, os conhecimentos tradicionais devem fazer parte do currículo, precisam ser valorizados e trabalhados de maneira transdisciplinar no currículo e, principalmente, precisam ser respeitados como conhecimentos válidos e importantes para a formação dos alunos.
Caso contrário, corremos o risco de consolidar visões históricas - representações estereotipadas.
A educação intercultural afeta não somente aos diferentes aspectos do currículo \textbf{explícito} – objetivos, conteúdos propostos, \textbf{métodos} e estilos de ensino, materiais didáticos utilizados, etc. –, como também o currículo oculto e as relações entre os diferentes agentes do processo educativo – professores/as, alunos/as, coordenadores/as, pais, agentes comunitários, etc.
Neste sentido, trabalhar os ritos, símbolos, imagens, etc., presentes no dia-a-dia da escola e a autoestima dos diferentes \textbf{sujeitos} e construir relações democráticas \textbf{que} superem o autoritarismo e o machismo tão fortemente arraigados nas culturas latino-americanas, constituem desafios iniludíveis ( CANDAU, 2008, p. 59 ).
Podemos destacar \textbf{que} avanços bem significativos no campo das políticas públicas, como estágio fundamental para a transformação das condições político-sociais, principalmente, quando relacionadas à garantia de direitos de minorias étnicas e/ou \textbf{identitárias}.
Nesse sentido, a Lei de Diretrizes e Bases da Educação Nacional ( Lei 9.394/96 ), as Diretrizes da Educação Escolar Indígena ( 2012 ), os Referenciais para Implantação de programas de Formação de Professores Indígenas nos Sistemas Estaduais de Ensino ( 2001 ), e legislações complementares asseguram o respeito e atendimento escolar \textbf{que}, dentre outros aspectos, fazem especial referência à inserção das línguas indígenas e valorização das mesmas, numa perspectiva de educação bilíngue. \textbf{Contudo}, a \textbf{tarefa} a ser realizada não é tão \textbf{simples}.
A complexidade das realidades \textbf{que} compõem os contextos da educação escolar indígena precisa, para \textbf{um} efetivo cumprimento dos objetivos de \textbf{uma} educação diferenciada e bilíngue, ser conhecidas.
O cenário brasileiro \textbf{atual} – com seus movimentos interculturais – tem \textbf{revelado} diferentes realidades sociolinguísticas, nas quais as línguas ancestrais se \textbf{configuram} tanto como primeira quanto segunda língua, além dos contextos em \textbf{que} estão caindo em desuso por seus povos, com \textbf{fortes} indícios de conflitos diglóssicos com a língua majoritária.
Neste contexto, atender a \textbf{atual} proposta da BNCC e das indicações das políticas públicas de educação em vigor \textbf{hoje} no Brasil, \textbf{implica} no reconhecimento do direito dos povos indígenas, como pessoas \textbf{que} usam \textbf{uma} língua diferente da língua majoritária, \textbf{que} têm valores/tradições/saberes singulares, e \textbf{que} têm o direito de serem \textbf{educadas} na sua língua, de terem sua identidade cultural incorporada aos conhecimentos e modos de funcionamento da

% matched lemmas: atravessar, atual, central, configurar, constitutivo, contudo, didático-pedagógico, educar, estímulo, et, explícito, forte, fronteira, hoje, identitária, implicar, inerente, instância, inter-relação, lado, lócus, método, neutro, permear, preciso, que, revelar, simples, sujeito, supor, tarefa, um, ver, vivar, viver
\end{textsample}
